% Configuração de layout correspondente ao modelo DOCX
\usepackage{newtxmath}
\usepackage{siunitx}
\usepackage{subcaption}
\usepackage{fontspec}
% No DOCX: Times New Roman (o LibreOffice do sistema usa Liberation Serif).
% Liberation Serif é uma substituição metricamente compatível disponível em Linux.
\IfFontExistsTF{Liberation Serif}{%
	\setmainfont{Liberation Serif}%
}{%
	\setmainfont{TeX Gyre Termes}%
}
\usepackage[portuguese]{babel}
\usepackage{geometry}
\geometry{
	a4paper,
	left=1.27cm,right=1.27cm,
	top=1.27cm,bottom=1.27cm,
	includehead,headheight=1.16cm,headsep=-0.17cm,
	footskip=0.75cm,
}
\usepackage{graphicx}
\usepackage{booktabs}
\usepackage{float} % [H]: figuras e tabelas ficam exatamente onde aparecem no código
\usepackage{xcolor}
\usepackage{fancyhdr}
\usepackage{titlesec}
\usepackage{enumitem}
\usepackage{amsmath}
\usepackage{microtype}
\microtypesetup{protrusion=false,expansion=false}
\usepackage[hidelinks]{hyperref}
\usepackage{cleveref}
\usepackage{xurl} % quebra de linha em URLs longos na bibliografia
\usepackage{tikz}
\usetikzlibrary{positioning,arrows.meta,calc}
\usepackage{amsmath,graphicx}
\usetikzlibrary{arrows.meta, positioning, shapes.geometric}
\usepackage{caption}

% Estilos dos diagramas de blocos (secoes/03_procedimentos.tex).
\tikzset{
  etapa/.style={draw, rounded corners=2pt, fill=blue!6, align=center,
                text width=5.6cm, minimum height=0.9cm, font=\small},
  io/.style={align=flush left, text width=5.2cm, font=\footnotesize},
  seta/.style={-{Stealth[length=2mm]}, thick},
  setaio/.style={-{Stealth[length=1.6mm]}, gray!70!black},
  titulo/.style={font=\footnotesize\bfseries, text=gray!60!black},
}

% Placeholders (remover antes de entregar): \tbd{texto} a vermelho no texto;
% \placeholder[largura]{altura}{descrição} no lugar de uma figura.
\newcommand{\tbd}[1]{\textcolor{red}{[#1]}}
\newcommand{\placeholder}[3][0.9\linewidth]{%
  \fcolorbox{red}{red!4}{\parbox[c][#2][c]{#1}{\centering\small\color{red}#3}}}

% Rascunhos de análise a rever: \rasc{texto}. Para passar tudo a preto,
% trocar "red" por "black" na linha seguinte.
\colorlet{rascunho}{red}
\newcommand{\rasc}[1]{{\color{rascunho}#1}}

% Lista compacta para as etapas que acompanham os diagramas de blocos.
\newlist{etapas}{enumerate}{1}
\setlist[etapas]{label=\textbf{\arabic*.}, leftmargin=*, labelsep=0.3em,
  topsep=0pt, itemsep=2pt, parsep=6pt,
  before={\setlength{\abovedisplayskip}{3pt}\setlength{\belowdisplayskip}{3pt}%
          \setlength{\abovedisplayshortskip}{0pt}\setlength{\belowdisplayshortskip}{3pt}}}

\setlength{\parindent}{0pt}
\setlength{\parskip}{6pt}
\linespread{1.0}
\hyphenpenalty=9000
\exhyphenpenalty=9000
\emergencystretch=3em
\frenchspacing

% Cabeçalho com texto, logótipo do DEM/Universidade de Aveiro, e traço horizontal.
\pagestyle{fancy}
\fancyhf{}
\fancyhead[L]{\raisebox{9pt}[0pt][0pt]{\fontsize{12}{14}\selectfont Licenciatura em Engenharia Aeroespacial. \textbf{Aerodinâmica}}}
\fancyhead[R]{\raisebox{8pt}[0pt][0pt]{\includegraphics[width=4.08cm]{imagens/ua_dem_cortado.png}}}
\fancyfoot[R]{\thepage\hspace*{0.50cm}}
\renewcommand{\headrulewidth}{0.5pt}
\renewcommand{\footrulewidth}{0pt}
\renewcommand{\headrule}{\vskip-8pt\hbox to\headwidth{\hspace*{-1.02cm}\rule{\dimexpr\headwidth+2.11cm\relax}{0.5pt}\hss}\vskip7.5pt}

% Níveis do documento original (sem numeração de secções).
\titleformat{\section}{\bfseries\fontsize{20}{23}\selectfont}{}{0pt}{}
\titlespacing*{\section}{0pt}{1pt}{1pt}
\titleformat{\subsection}{\bfseries\fontsize{14}{16}\selectfont}{}{0pt}{}
\titlespacing*{\subsection}{0pt}{1pt}{-3pt}

% Níveis de listas do DOCX: ponto cheio, seguido do marcador 'o'.
\setlist[itemize,1]{label=\textbullet,leftmargin=1.23cm,labelsep=0.36cm,labelwidth=0.25cm,topsep=1pt,partopsep=0pt,itemsep=5pt,parsep=0pt}
\setlist[itemize,2]{label={\normalfont o},leftmargin=1.25cm,labelsep=0.36cm,labelwidth=0.25cm,topsep=5pt,partopsep=0pt,itemsep=5pt,parsep=0pt}

% Título com sombreado cinzento e contorno preto, como no Word.
\newcommand{\tituloEmCaixa}[1]{%
	\par\vspace*{-1pt}% Ajustado para compensar o espaço do cabeçalho em LaTeX
	\noindent\makebox[\linewidth][c]{%
		{\setlength{\fboxsep}{2pt}\setlength{\fboxrule}{0.5pt}%
				\fcolorbox{black}{black!9}{%
					\parbox[c][0.96cm][c]{\dimexpr\linewidth+0.30cm-2\fboxsep-2\fboxrule\relax}{%
						\centering\bfseries\fontsize{25}{27}\selectfont #1}}}}%
	\par\vspace{-1pt}%
}

% O resumo é o único título de secção centrado no documento original.
\newcommand{\tituloResumo}{%
	\par\vspace{0pt}%
	\noindent\makebox[\linewidth][c]{\bfseries\fontsize{20}{23}\selectfont Resumo}%
	\par\vspace{-4pt}%
}
